\section{Conclusiones}

\noindent Se concluye que el Multi-Source Capacitated Facility Location Problem with Customer Incompatibilities (MS-CFLP-CI) corresponde a una variante relevante y compleja del problema de localización de instalaciones capacitadas. Su principal característica es permitir que la demanda de un cliente sea satisfecha por más de una bodega, incorporando además restricciones de incompatibilidad entre clientes. Esto permite representar escenarios logísticos más realistas, donde no solo se deben minimizar costos de apertura y operación, sino también respetar capacidades e impedir que ciertos pares de clientes sean atendidos por una misma instalación \cite{maia2023metaheuristic,ceschia2024multineighborhood}.\\

\noindent El estudio realizado permite distinguir la importancia de las variantes del problema. El CFLP clásico, el CFLP-CI, la variante single-source y la variante multi-source presentan diferencias relevantes en la forma de representar y resolver las asignaciones. En particular, la variante multi-source entrega mayor flexibilidad, pero también aumenta la complejidad, ya que las incompatibilidades deben verificarse en cada bodega donde exista asignación positiva de demanda.\\

\noindent La propuesta implementada, basada en una representación mediante matriz de envíos y en un algoritmo Hill Climbing con estrategia Best Improvement y reinicios, resultó adecuada para abordar el problema. La matriz de envíos permitió representar soluciones compatibles con la naturaleza multi-source, mientras que los operadores de búsqueda local permitieron modificar asignaciones, redistribuir demanda y mejorar progresivamente el costo total. Además, la construcción \(\epsilon\)-greedy permitió generar soluciones iniciales diversas sin abandonar completamente el criterio de bajo costo.\\

\noindent Los resultados experimentales justifican la relevancia de la propuesta. La solución greedy inicial permitió obtener soluciones factibles de manera rápida, pero su carácter determinístico limitó la exploración del espacio de búsqueda. La aplicación de Hill Climbing determinístico redujo el costo inicial, y la incorporación de reinicios estocásticos permitió obtener soluciones de menor costo en todas las instancias evaluadas. En las instancias pequeñas y medianas, la mejora promedio respecto de la solución greedy estuvo aproximadamente entre \(20.79\%\) y \(24.22\%\). En las instancias grandes, las mejoras fueron menores, entre \(11.18\%\) y \(14.74\%\), pero se mantuvieron consistentes pese al aumento del tamaño del vecindario.\\

\noindent También se concluye que el parámetro \(\epsilon\) influye directamente en la calidad de las soluciones obtenidas. Un valor nulo elimina la diversificación y genera siempre la misma solución inicial, mientras que un valor demasiado alto puede introducir aleatoriedad excesiva y deteriorar la calidad de las soluciones. En los experimentos realizados sobre wlp01.in, el valor \(\epsilon = 0.20\) obtuvo el mejor desempeño promedio, lo que evidencia un equilibrio adecuado entre intensificación greedy y diversificación estocástica.\\

\noindent El análisis de convergencia permitió observar que aumentar el número de iteraciones por reinicio mejora la calidad promedio de las soluciones, aunque con un mayor costo computacional. En la instancia wlp03.in, el costo promedio disminuyó progresivamente al aumentar las iteraciones, lo que confirma que la búsqueda local continúa encontrando mejoras durante el proceso. Sin embargo, el crecimiento del tiempo de ejecución en instancias grandes muestra que la exploración completa del vecindario puede volverse costosa.\\

\noindent Respecto del estado del arte, los métodos exactos, como la programación lineal entera mixta o la descomposición de Benders, son relevantes para formular el problema y obtener soluciones óptimas en instancias pequeñas o medianas, pero pueden presentar dificultades en instancias grandes o con restricciones adicionales \cite{fischetti2016benders}. Por ello, los enfoques heurísticos, metaheurísticos y matheurísticos resultan especialmente importantes para este tipo de problema, ya que permiten obtener soluciones de buena calidad en tiempos computacionales razonables \cite{maia2023metaheuristic,mansini2022kernel}.\\

\noindent Como trabajo futuro, sería relevante incorporar vecindarios más selectivos y estrategias metaheurísticas más robustas, como Simulated Annealing, Tabu Search o enfoques híbridos. También sería útil comparar distintas metaheurísticas bajo las mismas instancias, tiempos de ejecución y criterios de evaluación.